\clearpage
\phantomsection

\addcontentsline{toc}{chapter}{Tóm tắt}
\chapter*{\fontsize{13}{13}\selectfont{Tóm tắt}}

Trí tuệ nhân tạo (AI) đang trở thành xu hướng phát triển mạnh mẽ, trong đó việc ứng dụng mô hình Học tăng cường (Reinforcement Learning) để xây dựng các tác nhân thông minh (agents) đóng vai trò nền tảng quan trọng. So với các phương pháp lập trình kịch bản truyền thống (Scripted AI), các mô hình này khắc phục được nhiều hạn chế như hành vi rập khuôn, thiếu tính linh hoạt và khả năng thích ứng kém trước các tình huống bất ngờ. Thể loại game Roguelike với đặc tính tạo nội dung ngẫu nhiên và không gian trạng thái lớn cho phép xây dựng các agent có khả năng tự học, ra quyết định tối ưu và tổng quát hóa hiệu quả hơn. Trong báo cáo này, tác giả tập trung nghiên cứu cơ chế hoạt động của các thuật toán học tăng cường tiêu biểu (như PPO hoặc DQN), phân tích ưu điểm và thách thức khi triển khai cho các agent ảo. Trên cơ sở đó, báo cáo đề xuất và phát triển một mô hình thực nghiệm trong môi trường game Roguelike, qua đó minh chứng tiềm năng ứng dụng của Học tăng cường trong việc giải quyết các bài toán ra quyết định phức tạp.

\vspace{0.5cm}
\noindent{\textbf{Từ khóa:}} \textit{Học tăng cường (Reinforcement Learning), Game Roguelike, Procedural Content Generation (PCG), Deep Reinforcement Learning, PPO (Proximal Policy Optimization), Tác nhân tự chủ (Autonomous Agents)}